% DEFINED:   sec:nn-as-function, sec:linear-limits, eq:residual
% RESOLVED:  (none; first sections of the chapter)
% FORWARD:   sec:mlp-backprop (Task 4), sec:logistic (Task 5), sec:softmax (Task 5)

\section{What is a neural network?}\label{sec:nn-as-function}

A neural network is a parameterized function $f_\theta: \R^n \to \R^m$.
It takes a real-valued input vector and produces a real-valued output
vector. That is the whole structural commitment. What the output
\emph{means} is a separate decision, made by the loss.

The network produces raw values; the loss interprets those values as a
prediction by composing them with an \textbf{output activation} (also
called a \textbf{link function} in the language of generalized linear
models). The choice of link function and matching loss together specify
what kind of thing the network is predicting.

The simplest case is \textbf{regression}: the raw output \emph{is} the
prediction. The link function is the identity, and the loss is
\textbf{mean squared error}:
\begin{equation}\label{eq:residual}
  L = \frac{1}{2} \sum_i (\hat{y}_i - y_i)^2,
  \qquad
  \frac{\partial L}{\partial \hat{y}_i} = \hat{y}_i - y_i.
\end{equation}
The gradient is just the residual. This pattern, $\hat{y} - y$, recurs
throughout the chapter. It is the gradient of every canonical loss with
respect to its raw output, for reasons that turn out to be a theorem
about exponential families and canonical links, not an accident
(Chapter~2 proves it).

The \texttt{MSELoss} class encodes this directly:
\begin{lstlisting}
class MSELoss(Loss):
    def value(self, logits, y):
        return 0.5 * sum((zi - yi) ** 2 for zi, yi in zip(logits, y))

    def grad(self, logits, y):
        return [zi - yi for zi, yi in zip(logits, y)]

    def probs(self, logits):
        # Identity link: the prediction IS the raw output.
        return list(logits)
\end{lstlisting}
Note the $\frac{1}{2}$ convention in \texttt{value}: it makes the gradient
exactly $\hat{y} - y$ with no leading constant, matching the pattern of
the other losses.

A concrete demonstration: fit $y = \sin(x) + \varepsilon$ on
$x \in [-\pi, \pi]$ with a small MLP, identity output, and MSE loss
(\texttt{examples/regression.py}). After 2000 epochs of mini-batch SGD
the network learns the curve. This architecture is the same MLP we build
up to in \Cref{sec:mlp-backprop}; only the output head and loss differ.

With regression in hand as the unspecialized base case, the rest of the
chapter is a tour through two questions. First, what
\emph{expressive limits} does the architecture itself impose? That is
\Cref{sec:linear-limits,sec:mlp-backprop}: a single linear unit cannot
represent all continuous functions; the multilayer perceptron can.
Second, what \emph{output-side specializations} let the same
architecture solve concrete prediction problems beyond regression? That
is \Cref{sec:logistic,sec:softmax}: classification by sigmoid for
binary outcomes, softmax for categorical.

\section{The single linear unit and its limits}\label{sec:linear-limits}

The smallest non-trivial network is one \texttt{Linear} layer mapping
inputs to outputs:
\[
  \vz = W \vx + \mathbf{b}.
\]
This is a \textbf{single-layer perceptron} (SLP): an affine map with no
non-linearity. Trained with MSE it learns the least-squares fit. Trained
with a sigmoid head and BCE loss (\Cref{sec:logistic}), it is logistic
regression. The function class either way is \emph{linear} in the input.

The SLP can fit any linearly separable target. It cannot fit anything
else.

The canonical demonstration is XOR on $\{0, 1\}^2$:
\[
  \mathrm{XOR}(0,0) = 0, \quad
  \mathrm{XOR}(0,1) = 1, \quad
  \mathrm{XOR}(1,0) = 1, \quad
  \mathrm{XOR}(1,1) = 0.
\]
Training a single \texttt{Linear} layer with \texttt{SigmoidBCE} on these
four points for 500 epochs (batch size 4, seed~42):
\begin{lstlisting}
slp = Network([Linear(2, 1)], loss=SigmoidBCE())
slp.fit(xor_X, xor_Y, epochs=500, lr=0.5, batch_size=4)
\end{lstlisting}
Every input converges to $p = 0.5000$. The model literally cannot do
better. To see why, write down what fitting all four points would
require. With $\hat{y} = \sigma(w_1 x_1 + w_2 x_2 + b)$ and targets
$\{0, 1, 1, 0\}$, the BCE loss is minimized when each $\hat{y}$ matches
its target. Since sigmoid is monotone, the four constraints reduce to
four inequalities:
\[
  b < 0, \quad
  w_1 + b > 0, \quad
  w_2 + b > 0, \quad
  w_1 + w_2 + b < 0.
\]
Adding the second and third gives $w_1 + w_2 + 2b > 0$, so
$w_1 + w_2 + b > -b > 0$. But the fourth requires
$w_1 + w_2 + b < 0$. Contradiction. No choice of $(w_1, w_2, b)$
fits all four points simultaneously.

Geometrically: a linear classifier draws a single hyperplane in the
input space (a line in $\R^2$). XOR's positive class
$\{(0,1),(1,0)\}$ and negative class $\{(0,0),(1,1)\}$ lie at
diagonal corners of the unit square. No line separates them.

This was \textcite{minsky1969perceptrons}'s argument in 1969, which
stalled the perceptron research program for a decade and a half. The
resolution is one of the cleanest examples of how architecture and
expressivity interact. To escape linearity, add a non-linearity: the
subject of \Cref{sec:mlp-backprop}.
